% Transcription — fonds Alexandre Grothendieck, shelfmark 107, batch 1 (pages 1-20).
% Established from the facsimile archives/batches/107/batch-01.pdf, cut from
% a copy of 107.pdf fetched through the project's own relay
% (https://grothendieck-relay.onrender.com/source/107.pdf) — not uploaded by
% the user, and not mirrored with `npm run archive`, whose host is blocked —
% per the skill transcribe-grothendieck. The batch file carries no cover
% sheet: PDF page k is folder page k. The folder is twenty-five pages; this
% is the first of two batches (pages 21-25 are batch 2, transcribed in
% parallel and committed before this file was finished; its notation was
% read and followed: X (x) J, A x I, A^I, Delta(1), X_oo, ~, i a
% cofibration, p a fibration).
%
% Pass: Opus 5.5 (claude-opus-5-5), 2026-09-26 — first pass, unchecked against the pages by a human.
%
% The paper. Every leaf is continuous-feed computer paper from the CNUSC
% (Montpellier). Several listings carry a date: p. 1 « 06/04/82 » (the
% same job, PREDONTA, 15.27), p. 4 « 04 JUN 82 », p. 19 « VENDR. 04 JUIN
% 82 »; batch 2's p. 23 is stamped 82155, also 4 June 1982. The notes are
% therefore not earlier than June 1982, which fits the archivists'
% « [à partir de 1982] ». This is said once in a \note on page 1, not in
% \dating{}, which is copied verbatim from the catalogue.
%
% Pages skipped: 4, 6, 8, 10, 12 and 19 — the printed sides of JCL and
% job-log listings, with no mathematics of his (the pencil marks on them
% are corrections to the job control, « dummy », BLKSIZE values, and on
% p. 19 an isolated pencil Delta with an arrow). The gaps in the numbering
% are the record.
%
% Covers: page 1 (« Formalisme des X (x) M, X^M, Hom_A, Hom_B relatif : ...
% produits », ink over the header of a listing) and page 17 (« Question des
% lim filtrantes dans Hot », alone on a blank sheet) are headings for the
% runs that follow. They take no \page; their words become the two section
% headings, each with a \note. Page 1 is nonetheless where the listing's
% date is observed, and that note sits under the first heading.
%
% Pages transcribed: 2, 3, 5, 7, 9, 11, 13, 14, 15, 16, 18, 20.
% Pages summarised: none.
%
% His pagination. Two runs of circled or plain letters at the top of the
% leaves. The French run: a (p. 3), b (5), c (7), d (9), e (11), f (13),
% g (15) — the odd pages, the even ones being the printed backs. The
% English run: a (p. 18, circled), b (p. 20, circled); batch 2 carries e
% (p. 22) and d (p. 24), and no c has been found in either batch. Page 2
% (English, cancelled), 14 and 16 carry no letter. Each letter is given
% in a \note at its page.
%
% Language. Pages 3-15 are in French; pages 2, 18 and 20 are in English,
% as are batch 2's pages. His English is transcribed as it stands; notes
% and marginal apparatus are in French.
%
% What the batch is about. Two things. (1) Pages 2-16: the formalism of
% presheaves on a product A x B of small categories, set up for simplicial
% objects (B = Delta, « sA »): the projections p_A, p_B and the
% identifications (A x B)^ = Hom(A°, B^) = Hom(B°, A^); relative internal
% Homs Hom_A = p_B* Hom, Hom_B = p_A* Hom; the evaluation functors
% Gamma''_b and their extension Gamma''_M to presheaves M on B, with
% Gamma''_M(X) = p_B* Hom(p_A^*(M), X); the cotensor X^M = Hom(p_A^*(M), X)
% and, when A has sums, a tensor X (x) M with (X (x) M)(b) the coproduct
% over M(b) of copies of X(b), and the adjunction Hom_B(X, Y^M) = Hom_B(X
% (x) M, Y) for X, Y in Hom(B°, A); the external product L (box) M =
% p_B^*(L) x p_A^*(M), a Proposition that (L, M) |-> L (box) M is fully
% faithful, and the Corollaire Hom(L' (box) M', L (box) M) = Hom(L', L)
% (box) Hom(M', M), whence p_B^* commutes with internal Homs and L^M = L;
% a page on (X^{Delta_n})_0 = X_n and the coskeleton, with a question on
% Hom(M, a) => a; and a cancelled page on Hom_!(A^, M) (functors commuting
% with colimits) and its duals. Page 2, an English first start on the same
% formalism for A x Delta, is cancelled. (2) Pages 18-20 (English): do
% filtering (countable) colimits exist in (Hot), and do (Cat) -> Hot(W)
% and A^ -> Hot_A (A a test category) commute with them? Reduction to A =
% Delta through the nerve; replacing a sequence X_0 -> X_1 -> ... by
% monomorphisms; the limiting property (*) Hom_Hot(X_oo, Y) = lim
% Hom_Hot(X_i, Y); for Y a W-Kan object, injectivity of (*) reduces to
% rectifying a square i : X -> X' (cofibration), p : Y -> Y' (fibration)
% commuting up to homotopies k, h — « exactly the content of Quillen's
% version of "covering homotopy extension theorems" for simplicial closed
% model categories (prop. 4 of Chap. II, n° 2.4) ». Batch 2 continues this
% square (p. 22, e) and the injectivity (p. 24, d).
%
% On batch 2. Page 24 (d) opens « Returning to our initial pb » and shows
% Hom(X_oo, Y) -> lim Hom(X_i, Y) « surjection »: that is (*) of p. 18 and
% the « Injectivity of (*)? » of p. 20 here. The square of p. 20 (g, phi,
% i, p, k, h) is the one batch 2 draws on pp. 22-23. Nothing here
% contradicts batch 2's readings; batch 2 was not edited.
%
% Notation fixed once for the batch: A^ as $A^{\wedge}$ (he writes a
% circumflex after the letter; $\hat{B}$ once, p. 5, kept); A° as
% $A^{\circ}$; internal Hom underlined, $\underline{\mathrm{Hom}}$; his
% boxed cross as $\boxtimes$; (Ens), (Cat), (Hot) with parentheses as he
% writes them; Hot_A and (Hot)_A both as written; W underlined as
% $\underline{W}$; the simplex category as $\Delta$ (he draws it with a
% doubled stroke), the boundary as $\dot{\Delta}_n$. Greek letters are in
% math mode throughout.
%
% Readings to check first: the cover of p. 1 (third line); p. 3 « On ...
% les diagrammes ci-dessus » (three words \ill{}); p. 5, the rewritten
% first line; p. 7, the opening line and the footnote marked *; p. 9, the
% interlinear words before « Y^M, X (x) M »; p. 11, which of p_A*, p_B*
% goes in the formulas for gX(a), dX(b); p. 13, the marginal note on ! and
% !!; p. 15, the words around « flèche canonique »; p. 16, the exponent
% « 1 » on A; p. 18, the sentence on (ii); p. 20, the left-margin note and
% the l / r over the two ~.

\documentclass[11pt,a4paper]{article}
% Le préambule commun à toutes les transcriptions du fonds Grothendieck.
%
% Il tient en une idée : **l'incertitude se compose, elle ne se résout pas.**
% Une transcription qui lisse un mot douteux a détruit la seule chose pour
% laquelle on la faisait — un lecteur doit pouvoir savoir, à chaque endroit,
% ce qui était sur le feuillet et ce qui a été ajouté. Les six macros qui
% suivent sont donc l'appareil critique, et non de la décoration.
%
% Les mêmes commandes sont reconnues par `scripts/render.mjs`, qui produit la
% vue de lecture du site. Une macro ajoutée ici doit l'être aussi là-bas,
% faute de quoi le rendu échouera bruyamment — ce qui est voulu : un rendu
% silencieusement faux est pire qu'un rendu absent.

\ProvidesPackage{grothendieck}[2026/08/08 Transcriptions du fonds Grothendieck]

\usepackage{amsmath,amssymb,amsthm}
\usepackage{mathrsfs}
\usepackage{stmaryrd}
% Les diagrammes commutatifs sont partout dans ce fonds : c'est la langue
% même de la géométrie algébrique de Grothendieck, et une transcription qui
% les rendrait en prose ne transcrirait rien.
\usepackage{tikz-cd}
% Français par défaut — c'est la langue des feuillets — mais l'anglais est
% chargé aussi : la traduction et le résumé sont des documents anglais, et la
% typographie française (espaces avant les deux-points) y serait une faute.
% Ces documents basculent par \selectlanguage{english} en tête de corps.
\usepackage[english,french]{babel}
\usepackage{fontspec}
\usepackage{geometry}
\geometry{a4paper,margin=25mm,marginparwidth=28mm}
\usepackage{marginnote}
\usepackage{xcolor}
% ulem, not soul. `soul`'s \st and \ul are built on a character-by-character
% scan that XeLaTeX's font machinery breaks: the document fails with an
% undefined control sequence rather than degrading. ulem does the same two jobs
% and is Unicode-engine safe. `normalem` stops it hijacking \emph, which the
% transcriptions use for Grothendieck's own emphasis.
\usepackage[normalem]{ulem}
\usepackage{hyperref}
\hypersetup{colorlinks=true,linkcolor=black,urlcolor=black,pdfborder={0 0 0}}

% --- Métadonnées du lot -----------------------------------------------------
% Reprises telles quelles par le script de rendu, qui les affiche en tête de
% la vue de lecture. Elles fixent ce que le fichier prétend couvrir : sans
% elles, une transcription flotte sans point d'attache dans un fonds de seize
% mille feuillets.

\newcommand{\folder}[1]{\gdef\grothFolder{#1}}
\newcommand{\batch}[1]{\gdef\grothBatch{#1}}
\newcommand{\leaves}[2]{\gdef\grothFirst{#1}\gdef\grothLast{#2}}
\newcommand{\foldertitle}[1]{\gdef\grothFoldertitle{#1}}
\gdef\grothFolder{?}\gdef\grothBatch{?}\gdef\grothFirst{?}\gdef\grothLast{?}\gdef\grothFoldertitle{}

% --- Le résumé --------------------------------------------------------------
%
% Il ouvre la lecture modernisée, sous le titre « Résumé », et il est composé
% en petit. Ce n'est pas un ornement : c'est le seul passage du document qui
% ne soit pas des mathématiques, et un lecteur doit pouvoir le reconnaître
% avant de l'avoir lu — puis le sauter s'il connaît déjà le sujet, ou s'y
% arrêter s'il ne le connaît pas. Le filet qui le suit marque la reprise du
% corps.

\newenvironment{resume}
  {\small\setlength{\parskip}{0.45em}}
  {\par\medskip\hrule\medskip}

% --- Mots-clés ---------------------------------------------------------------
%
% En anglais, en clôture du résumé de la lecture modernisée : le vocabulaire
% moderne sous lequel chercher ce que les feuillets construisent. Le manifeste
% du site (`npm run manifest`) les extrait pour étiqueter la cote entière —
% cette ligne est la source unique des tags, il n'y a pas de fichier à côté.
\newcommand{\keywords}[1]{\par\smallskip\noindent
  {\footnotesize\textsc{Keywords} --- \textit{#1}}\par}

% --- L'appareil critique ----------------------------------------------------

% Le numéro de feuillet, en marge. C'est l'ancre qui permet de régler un
% désaccord sur une formule en regardant *un* feuillet plutôt qu'une chemise
% de six cents. Le site s'en sert aussi pour tourner les pages du fac-similé
% quand on fait défiler la transcription.
\newcommand{\leaf}[1]{%
  \marginnote{\footnotesize\color{gray!70}\textsf{#1}}%
  \label{leaf:#1}\ignorespaces}

% Illisible : on ne devine pas. Un mot inventé qui se lit comme les autres est
% le pire résultat possible.
\newcommand{\ill}{\textcolor{red!60!black}{[\,\ldots\,]}}

% Lecture douteuse : le mot est donné, et signalé comme incertain.
\newcommand{\uncertain}[1]{\uline{#1}}

% Ajout de l'éditeur — une lettre manquante, un mot sous-entendu.
\newcommand{\add}[1]{\textcolor{blue!50!black}{[#1]}}

% Biffé par Grothendieck lui-même. Ce qu'il a rayé fait partie du document :
% une rature dit souvent où la pensée a changé de direction.
\newcommand{\struck}[1]{\sout{#1}}

% Note du transcripteur, sur l'état matériel du feuillet.
\newcommand{\note}[1]{\par\smallskip{\small\color{gray!60!black}\textit{[#1]}}\par\smallskip}

% Note marginale de Grothendieck, distincte du corps du texte.
\newcommand{\marginal}[1]{\par\smallskip\noindent\hspace{1em}%
  {\small\color{gray!60!black}#1}\par\smallskip}

% --- Titre ------------------------------------------------------------------

\AtBeginDocument{%
  \begin{center}
    {\large\bfseries Fonds Alexandre Grothendieck --- cote n\textsuperscript{o}~\grothFolder}\\[2pt]
    {\itshape\grothFoldertitle}\\[4pt]
    {\small feuillets \grothFirst--\grothLast\ (lot \grothBatch)}\\[2pt]
    {\footnotesize\color{gray!60!black}Transcription établie d'après le fac-similé numérisé
     par l'Université de Montpellier.\\
     Les lectures incertaines sont soulignées, les illisibles notées [\,\ldots\,],
     les ajouts entre crochets.}
  \end{center}
  \vspace{1em}}

\endinput


\folder{107}
\batch{1}
\pages{1}{20}
\dating{[à partir de 1982]}
\watermark{Édition de démonstration}
\foldertitle{Closed models [notes postérieures à PS ?] : notes manuscrites (s.d.)}

\begin{document}

\section{Formalisme des $X \otimes M$, $X^{M}$, $\underline{\mathrm{Hom}}_A$, $\underline{\mathrm{Hom}}_B$ relatif}

\note{titre de sa main, à l'encre, sur la page 1, écrit par-dessus
l'en-tête d'un listing ; il continue : « relatif : \ill{}
\uncertain{définitions} » et, souligné, « \uncertain{produits} ». Le
listing de cette page est daté « 06/04/82 » ; celui de la page 19,
du même travail, « 04 JUIN 82 » : les feuilles de ce lot ne sont pas
antérieures à juin 1982}

\page{2}

\note{toute la page est annulée par deux longs traits obliques ; elle est
transcrite telle qu'elle se lit. Elle est en anglais}

\begin{tikzcd}[column sep=small]
 & A \times \Delta \arrow[dl, "p_{\Delta}"'] \arrow[dr, "p_A"] & \\
A \arrow[dr, "q_A"'] & & \Delta \arrow[dl, "q_{\Delta}"] \\
 & e &
\end{tikzcd}

\note{le sommet de gauche est de lecture douteuse : \uncertain{$A$}}

\[
X_*, Y_* \in \mathrm{Ob}\, \underbrace{(A \times \Delta)^{\wedge}}
\simeq \underline{\mathrm{Hom}}(\Delta^{\circ}, A^{\wedge})
\simeq \underline{\mathrm{Hom}}(A^{\circ}, \Delta^{\wedge})
\]
\[
\begin{cases}
\mathrm{Hom}_{A \times \Delta}(X_*, Y_*) = \Gamma_{A \times \Delta}\bigl(\underline{\mathrm{Hom}}_{A \times \Delta}(X_*, Y_*)\bigr)\\
\underline{\mathrm{Hom}}_A(X_*, Y_*) = p_{\Delta*}\, \underline{\mathrm{Hom}}_{A \times \Delta}(X_*; Y_*)\\
\underline{\mathrm{Hom}}_{\Delta}(X_*, Y_*) = p_{A*}\, \underline{\mathrm{Hom}}_{A \times \Delta}(X_*, Y_*)
\end{cases}
\]

\[
\begin{cases}
p_{A*} X_* = (\Gamma_A X_n)_n\\
p_{\Delta*} X_* = X_0
\end{cases}
\qquad
(q_{\Delta} p_A)_*(X_*) = (q_A p_{\Delta})_* X_* = \Gamma_A X_0
\]
\note{ces formules, en haut à droite, sont isolées du reste par un trait
anguleux ; dans la première, le $X_*$ est repassé d'une rature, et le
dernier indice, au bord de la feuille, est de lecture douteuse}

\begin{tikzcd}[column sep=small]
 & (A \times \Delta)^{\wedge} \supset s\underline{A} = \underline{\mathrm{Hom}}(\Delta^{\circ}, A) \arrow[dl, "p_{\Delta*}"'] \arrow[dr, "p_{A*}"] & \\
A^{\wedge} \arrow[dr, "q_{A*}"'] & & \Delta^{\wedge} \arrow[dl, "q_{\Delta*}"] \\
 & (\mathrm{Ens}) &
\end{tikzcd}

\note{devant $(A \times \Delta)^{\wedge}$, un mot raturé, \struck{\ill{}}}

\struck{\ill{}} What about $\bigl(\underline{\mathrm{Hom}}_A(L, X_n)\bigr)_n
\overset{?}{\in} (A \times \Delta)^{\wedge}$,
$\bigl(\mathrm{Hom}_A(L, X_n)\bigr)_n \in \Delta^{\wedge}$ for given
$L \in \mathrm{Ob}\, A^{\wedge}$ (more generally, $L \in \mathrm{Ob}\,
A^{\wedge}$ ?)\note{les exposants des deux $A$ sont surchargés ; le second
est suivi d'un trait oblique et d'un point d'interrogation}: the second is
the transform of the first by $p_{A*}$ (${}= \Gamma_A$ applied
componentwise). I contend that if we identify $L$ with the constant complex
with value $L$, i.e.\ with $p_{\Delta}^{*}(L)$, we get
\[
\begin{aligned}
&\text{I}\quad \bigl(\underline{\mathrm{Hom}}_A(L, X_n)\bigr)_{n \in \mathbb{N}}
= \underline{\mathrm{Hom}}_{A \times \Delta}\bigl(p_{\Delta}^{*}(L), X_*\bigr)
&&\text{hence, by applying } p_{A*}\\
&\text{II}\quad \bigl(\mathrm{Hom}_A(L, X_n)\bigr)_{n \in \mathbb{N}}
= \underline{\mathrm{Hom}}_{\Delta}\bigl(p_{\Delta}^{*}(L), X_*\bigr)
\end{aligned}
\]
\note{devant $p_{A*}$, un premier indice biffé, \struck{$\Gamma_{A \times
\Delta}$} (lecture douteuse)}

To get the first one, we apply the definition of
$\underline{\mathrm{Hom}}$, we have to get the $n$-th component of the
second hand side; \struck{\ill{}} This is \uncertain{justified} easily
checked (see below) to be\note{la phrase s'arrête là. « justified » est
ajouté au-dessus de la ligne et relié par une boucle}

\note{dans la marge gauche, écrits en travers (la feuille tournée d'un
quart de tour), des fragments sans lien apparent avec le texte : « A. »,
\uncertain{« i.e. »}, « I », $\sigma^{\circ} \longleftrightarrow \sigma$ ;
$M(b)$ ; $X(a, b) = M(b)$ ; une ébauche de flèches
$\{i_0\} \to I$, $\underline{\mathrm{Hom}}(a_1, a_0) \to$ \ill{},
$u \mapsto u(i_0)$, et un signe raturé}

\page{3}

\note{en haut, au milieu, la lettre « a »}

\struck{$(A \times B)^{\wedge} \simeq \underline{\mathrm{Hom}}(A^{\circ},
B^{\wedge}) \simeq \underline{\mathrm{Hom}}(B^{\circ}, A^{\wedge})$}.

\begin{tikzcd}[column sep=small]
 & A \times B \arrow[dl, "p_B"'] \arrow[dr, "p_A"] & \\
A \arrow[dr, "q_A"'] & & B \arrow[dl, "q_B"] \\
 & e &
\end{tikzcd}

\begin{tikzcd}[column sep=tiny, row sep=large, nodes={font=\small}]
\underline{\mathrm{Hom}}(A^{\circ}, B^{\wedge}) \simeq \arrow[d, "{X' \mapsto \Gamma_B \circ X'}"'] & (A \times B)^{\wedge} \arrow[dl, "(p_B)_*"] \arrow[dr, "(p_A)_*"'] & \simeq \underline{\mathrm{Hom}}(B^{\circ}, A^{\wedge}) \arrow[d, "{X'' \mapsto \Gamma_A \circ X''}"] \\
\underline{\mathrm{Hom}}(A^{\circ}, (\mathrm{Ens})) \simeq A^{\wedge} \arrow[dr, "{q_{A*} = \Gamma_A}"'] & & B^{\wedge} = \underline{\mathrm{Hom}}(B^{\circ}, (\mathrm{Ens})) \arrow[dl, "{q_{B*} = \Gamma_B}"] \\
 & (\mathrm{Ens}) &
\end{tikzcd}

\note{à droite de la ligne du haut, entre crochets :
$\longrightarrow \underline{\mathrm{Hom}}(B^{\circ}, A)$, la flèche coupée
d'un petit trait vertical à crochets, de sens incertain. Sous les deux
flèches du bas : $q_{A*} = \Gamma_A \simeq
\underline{\mathrm{Hom}}_A(e_{A^{\wedge}}, \cdot)$, devant lequel une
lettre biffée, et $q_{B*} = \Gamma_B \simeq
\underline{\mathrm{Hom}}_B(e_{B^{\wedge}}, \cdot)$}

\[
\begin{aligned}
\underline{\mathrm{Hom}}_B(X, Y) &\overset{\text{déf}}{=} p_{A*}\, \underline{\mathrm{Hom}}_{A \times B}(X, Y)\\
\underline{\mathrm{Hom}}_A(X, Y) &\overset{\text{déf}}{=} p_{B*}\, \underline{\mathrm{Hom}}_{A \times B}(X, Y)
\end{aligned}
\]
\struck{$\mathrm{Hom}(X, Y) =$} \uncertain{donc}
\[
\mathrm{Hom}(X, Y) \simeq \Gamma_A\, \underline{\mathrm{Hom}}_A(X, Y)
\simeq \Gamma_B\, \underline{\mathrm{Hom}}_B(X, Y)
\]
\note{dans la dernière formule, l'indice du premier $\Gamma$ est surchargé ;
on lit $A$}

$A$ et $B$ deux petites catégories, on s'intéresse \uncertain{surtout}
\uncertain{aux} \uncertain{deux} $(A \times B)^{\wedge} \simeq
\underline{\mathrm{Hom}}(A^{\circ}, B^{\wedge}) \simeq
\underline{\mathrm{Hom}}(B^{\circ}, A^{\wedge})$. On \ill{} \ill{} \ill{}
les diagrammes ci-dessus \uncertain{quelques} propriétés des foncteurs
\ill{} divers. Pour $b \in B$, on a un foncteur
\[
\begin{array}{ccc}
\Gamma''_b : (A \times B)^{\wedge} & \longrightarrow & A^{\wedge}\\
\downarrow & & \|\\
\simeq \underline{\mathrm{Hom}}(B^{\circ}, A^{\wedge}) & \longrightarrow & A^{\wedge}\\
X'' & \longmapsto & X''(b)
\end{array}
\qquad \text{déduit par } (A \times B)^{\wedge} \simeq
\underline{\mathrm{Hom}}(A^{\circ}, B)
\]
\note{devant $\Gamma''_b$, un signe biffé ; sur la flèche du haut, un indice biffé ; à droite de $X''(b)$, un
mot biffé, \struck{\ill{}}. Dans la dernière formule, on lit $B$ et non
$B^{\wedge}$}

Je dis que l'on a
\[
\Gamma''_b(X) \simeq p_{B*}\, \underline{\mathrm{Hom}}\bigl(p_A^{*}(b), X\bigr)
\qquad \text{fonctoriel en } b, X
\]
\[
\bigl(\simeq \underline{\mathrm{Hom}}_A(p_A^{*}(b), X)\bigr),
\]
ce qui signifie aussi que l'on a
\[
\begin{aligned}
\underset{X(a, b)}{\underset{\|}{\mathrm{Hom}_{A^{\wedge}}\bigl(a, \Gamma''_b(X)\bigr)}}
&\simeq \mathrm{Hom}_{A^{\wedge}}\bigl(a, p_{B*}\, \underline{\mathrm{Hom}}(p_A^{*}(b), X)\bigr)
&&\text{fonctoriel en } a, X, b\\
&\simeq \mathrm{Hom}_{(A \times B)^{\wedge}}\bigl(p_B^{*}(a), \underline{\mathrm{Hom}}(p_A^{*}(b), X)\bigr)\\
&\simeq \mathrm{Hom}_{(A \times B)^{\wedge}}\bigl(p_B^{*}(a) \times p_A^{*}(b), X\bigr)
&&\text{OK}
\end{aligned}
\]
\note{plusieurs indices $A$, $B$ de ces lignes sont surchargés, l'un
récrit sur l'autre ; on donne la lecture finale}

\note{le paragraphe qui suit est annulé par trois traits obliques ; il est
transcrit tel qu'il se lit}

Symétriquement, on a des $\Gamma'_a$, $a \in A$
\[
\begin{array}{ccc}
\Gamma'_a : (A \times B)^{\wedge} & \longrightarrow & B^{\wedge}\\
\wr & &\\
\underline{\mathrm{Hom}}(A^{\circ}, B^{\wedge}) & &\\
X' & \longmapsto & X'(a)
\end{array}
\]
et \uncertain{aussi}
\[
\Gamma'_a(X) \simeq p_{A*}\bigl(\underline{\mathrm{Hom}}(p_B^{*}(a), X)\bigr)
\]

\page{5}

\note{en haut, au milieu du bord perforé, la lettre « b »}

\struck{On \ill{} \ill{} \ill{} \ill{} \ill{} \ill{} \ill{}}
\uncertain{On a une troisième interprétation du foncteur} $\Gamma''_b$, en
\struck{$\underline{\mathrm{Hom}}_{(A \times B)^{\wedge}}\bigl(p_A^{*}(b),
\underline{\mathrm{Hom}}(p_A^{*}(b), X)\bigr)$} suivant
$(A \times B)^{\wedge} \simeq \underline{\mathrm{Hom}}(A^{\circ},
B^{\wedge})$ et considérant \struck{\ill{}}
\[
\varepsilon_b : B^{\wedge} \to \mathrm{Ens}, \qquad
\varepsilon_b(G) = \mathrm{Hom}(b, G) = G(b)
\]
\note{la première ligne est biffée et récrite au-dessus ; la nouvelle
lecture est douteuse}

alors $\Gamma''_b(X')$ \struck{$(a)$} $= \varepsilon_b \circ X'$.

Ceci se généralise au cas où on se donne $M \in \mathrm{Ob}\, \hat{B}$,
quand on considère
\[
\varepsilon_M : B^{\wedge} \to \mathrm{Ens}, \qquad
G \mapsto \varepsilon_M(G) = \mathrm{Hom}(M, G)
\]
considérons donc, pour $X'' : A^{\circ} \to B^{\wedge}$ \uncertain{variable}
\note{sous $B^{\wedge}$, un astérisque}
\[
\Gamma''_M : \underset{\underline{\mathrm{Hom}}(A^{\circ}, B^{\wedge})}{\underset{\cap}{X''}}
\longmapsto \varepsilon_M \circ \Gamma''_M
\]
\note{on lit bien $\varepsilon_M \circ \Gamma''_M$ à droite, et non
$\varepsilon_M \circ X''$}

Je dis qu'on a, \uncertain{canoniquement}, $X \in \mathrm{Ob}\,(A \times
B)^{\wedge}$ \uncertain{variable}
\[
\Gamma''_M(X) \simeq p_{B*}\, \underline{\mathrm{Hom}}\bigl(p_A^{*}(M),
X\bigr) \qquad \text{fonctoriel en } M, X
\]
Pour le voir, \struck{considérons} il faut \uncertain{trouver}
\[
\begin{array}{ccc}
\mathrm{Hom}_{A^{\wedge}}\bigl(a, \Gamma''_M(X)\bigr) & \simeq &
\mathrm{Hom}_{A^{\wedge}}\bigl(a, p_{B*}\, \underline{\mathrm{Hom}}(p_A^{*}(M), X)\bigr)\\
\| & & \wr\\
\Gamma''_M(X)(a) & &
\mathrm{Hom}_{(A \times B)^{\wedge}}\bigl(p_B^{*}(a), \underline{\mathrm{Hom}}(p_A^{*}(M), X)\bigr)\\
\| & & \wr\\
\mathrm{Hom}_{B^{\wedge}}\bigl(M, X''(a)\bigr) & &
\mathrm{Hom}_{(A \times B)^{\wedge}}\bigl(p_B^{*}(a) \times p_A^{*}(M), X\bigr)\\
\wr & &\\
p_{A*}\bigl(\underline{\mathrm{Hom}}(p_B^{*}(a), X)\bigr) & &\\
\| & &\\
\mathrm{Hom}_{(A \times B)^{\wedge}}\bigl(p_A^{*}(M), \underline{\mathrm{Hom}}(p_B^{*}(a), X)\bigr) & &\\
\| & &\\
\mathrm{Hom}_{(A \times B)^{\wedge}}\bigl(p_A^{*}(M) \times p_B^{*}(a), X\bigr) & &
\end{array}
\]
fonct.\ en $M, X, a$\note{ces mots sont écrits à droite de la première ligne}
\note{colonne de gauche : devant $\Gamma''_M(X)(a)$, un « Hom( » biffé ;
dans $\mathrm{Hom}_{B^{\wedge}}(M, X''(a))$, le $X$ est surchargé, et sous
$X''(a)$ une flèche $\longmapsto$ porte un signe $=$ ; le terme
$p_{A*}(\ldots)$ est raturé et de lecture douteuse. Une double flèche
courbe, marquée « OK ! », relie la dernière ligne de gauche à la dernière
de droite}

\page{7}

\note{en haut, au milieu du bord perforé, la lettre « c »}

Nous nous intéresserons \uncertain{maintenant} à des propriétés de
stabilité \struck{\ill{}} \uncertain{des} sous-catégories
$\underline{\mathrm{Hom}}(B^{\circ}, A)$ de $(A \times B)^{\wedge}$, et $A$
de $A^{\wedge}$.\note{au-dessus du mot biffé, un mot ajouté, lu
\uncertain{pleines}} Le foncteur $p_{B*}$ $\bigl(= (\Gamma_B)^{\circ}\bigr)$
\ill{} \struck{$(A \times B)^{\wedge}$} $\underline{\mathrm{Hom}}(B^{\circ},
A)$ dans $A$, dès que $B$ admet un objet final $e_B$ (car il s'identifie
\uncertain{alors} à $\Gamma''_{e_B}$). Donc on trouve

\begin{tikzcd}[column sep=small]
 & \underline{\mathrm{Hom}}(B^{\circ}, A) \arrow[dl] \arrow[dr, "(\Gamma_A)^{\circ}"] & \\
A \arrow[dr, "\Gamma_A"'] & & B^{\wedge} = \underline{\mathrm{Hom}}(B^{\circ}, (\mathrm{Ens})) \arrow[dl, "\Gamma_B"] \\
 & \phantom{(\mathrm{Ens})} &
\end{tikzcd}

\note{le sommet de gauche est un $A$ dont l'exposant a été noirci d'une
rature — sans doute un $A^{\wedge}$ corrigé en $A$. Les deux flèches du bas
convergent vers un point où aucun terme n'est écrit}

\underline{NB} On pense au cas où $B = \Delta$ p.\ ex.\ (cat.\ des
simplexes-types) d'où $\underline{\mathrm{Hom}}(B^{\circ}, A) = sA$ objets
simpliciaux de $A$.

\struck{Je \ill{} \ill{}}

\struck{On} On s'intéressera \struck{\ill{}} au cas où
$\underline{\mathrm{Hom}}(B^{\circ}, A)$ est stable par
\[
\underline{\mathrm{Hom}}(p_A^{*}(M), X) \overset{\text{déf}}{=} X^{M}
\qquad \text{pour } \forall\, M \in \mathrm{Ob}\, B^{\wedge}
\]
(moyennant év\add{entuellemen}t des conditions de finitude) (alors que
$\underline{\mathrm{Hom}}(B^{\circ}, A)$ ne contient pas nécessairement les
$p_A^{*}(M)$ \ldots${}^{*}$).

Il suffit que $A$ soit stable par \uncertain{limites} du type requis
\add{…} pour représenter $M$ (en termes de foncteurs représentables).
\note{sous la ligne, un mot \ill{} suivi de $\varprojlim$ souligné, relié
par un trait courbe à la ligne suivante, et une flèche $\longleftarrow$
sous « limites » ; un double trait vertical marque ces deux lignes dans la
marge gauche}

Notons que \uncertain{l'on a}
\[
\begin{aligned}
\underline{\mathrm{Hom}}_B(X, Y^{M})
&\simeq p_{A*}\, \underbrace{\underline{\mathrm{Hom}}\bigl(X, \underline{\mathrm{Hom}}(p_A^{*}(M), Y)\bigr)}\\
&\simeq p_{A*}\bigl(\underline{\mathrm{Hom}}(X \times p_A^{*}(M), Y)\bigr)\\
&\simeq \underline{\mathrm{Hom}}_B(X \times p_A^{*}(M), Y)
\end{aligned}
\]

${}^{*}$ Pour contenir les $p_A^{*}(M)$, il faudrait que $\forall\, b$
\ill{} (\uncertain{quelconque}), le foncteur \underline{constant} sur $A$
de valeur $M(b)$ soit représentable. Mais cela signifie que $M(b)$ est
\ill{} \ill{} et $A$ \uncertain{a} objet final. \uncertain{Donc} \ill{} $M$
objet final de $B^{\wedge}$, et $A$ admet objet final !!\note{note de bas
de page, serrée sur trois lignes contre le bord perforé}

\page{9}

\note{en haut, au milieu du bord perforé, la lettre « d »}

\uncertain{donc}
\[
\underline{\mathrm{Hom}}_B(X, Y^{M}) \simeq
\underline{\mathrm{Hom}}_B(\underbrace{X \times p_A^{*}(M)}, Y)
\]
\note{sous l'accolade, une notation raturée à l'encre épaisse, où se
devine \struck{$X \otimes M$}}

Formule valable pour $\forall\, X, Y \in (A \times B)^{\wedge}$, mais
\uncertain{si}, \uncertain{comme} nous nous intéressons à cette formule pour
$X, Y \in \underline{\mathrm{Hom}}(B^{\circ}, A)$, \ill{} \ill{} \ill{}
\uncertain{y trouver} \uncertain{aussi} $Y^{M}$, $X \otimes M$ dans
$\underline{\mathrm{Hom}}(B^{\circ}, A)$. On a vu des conditions
\uncertain{naturelles} pour qu'il en soit ainsi pour $Y^{M}$ (existence de
certaines $\varprojlim$ dans $A$, à l'exclusion des
$\underline{\mathrm{Hom}}$ internes dans $A$), il reste à regarder pour
$X \times p_A^{*}(M)$, qui \uncertain{en tant que} foncteur $B^{\circ} \to
A^{\wedge}$ s'identifie à
\[
b \longmapsto \underbrace{X(b) \times M(b)} = \Bigl\{ a \longmapsto
\underbrace{\mathrm{Hom}_A\bigl(a, X(b)\bigr)} \times M(b) \Bigr\}
\]
\[
X(b)(a) = X(a, b)
\]
\note{la dernière égalité est écrite sous l'accolade de droite, qui porte
sur $\mathrm{Hom}_A(a, X(b))$ ; devant ce $\mathrm{Hom}$, un $X$ biffé}

Comme \uncertain{tel}, il n'est pratiquement jamais représentable. Mais
supposons que dans $A$ les sommes existent (\uncertain{sommes} finies
suffisent si les $M(b)$ sont \uncertain{à valeurs} finies) et définissons
\[
X \otimes M \in \underline{\mathrm{Hom}}(B^{\circ}, A) \quad \text{par}
\]
\[
(X \otimes M)(b) = \coprod^{A}_{\sigma \in M(b)} X(b)
\quad \Bigl( \xrightarrow[\;i_b\;]{\;\text{nat}\;}
\coprod^{A^{\wedge}}_{\sigma \in M(b)} X(b) = \bigl(X \times p_A^{*}(M)\bigr)(b) \Bigr)
\]
\note{le membre de gauche est surchargé, et un premier terme du membre de droite est biffé ; l'indice sous la flèche est lu
\uncertain{$i_b$}. Les deux sommes portent en exposant la catégorie où
elles sont prises, $A$ et $A^{\wedge}$}

alors on a, pour $X, Y$ dans $\underline{\mathrm{Hom}}(B^{\circ}, A)$
\[
\underline{\mathrm{Hom}}_B(X, Y^{M}) \simeq \underline{\mathrm{Hom}}_B(X
\otimes M, Y) \qquad \text{fonctoriel}
\]

\page{11}

\note{feuillet tenu en largeur ; tout est écrit sur la moitié droite. En
haut, au milieu, la lettre « e »}

\textbf{Récapitulation}\note{le titre, souligné, est en haut à droite}

\begin{tikzcd}
 & (A \times B)^{\wedge} \arrow[dl, "p_{B*}"'] \arrow[dr, "p_{A*}"] & \\
A^{\wedge} \arrow[ur, bend left=50, "p_B^{*}"] \arrow[dr, "q_{A*} = \Gamma_A"'] & & B^{\wedge} \arrow[ul, bend right=50, "p_A^{*}"'] \arrow[dl, "q_{B*} = \Gamma_B"] \\
 & (\mathrm{Ens}) &
\end{tikzcd}

\[
\begin{aligned}
X(a, b) &= \Gamma_{A \times B}\,
\underline{\mathrm{Hom}}\bigl(p_B^{*}(a) \times p_A^{*}(b), X\bigr)\\
&= \Gamma_{A \times B}\bigl(\underline{\mathrm{Hom}}_{A \times B}(a \boxtimes b, X)\bigr)
\end{aligned}
\qquad
\begin{cases}
= \Gamma_A\, p_{B*}\, \underline{\mathrm{Hom}} \\
= \Gamma_B\, p_{A*}\, \underline{\mathrm{Hom}}
\end{cases}
\]
\note{au-dessus de $p_B^{*}(a)$ et de $p_A^{*}(b)$, deux accolades portant
« \uncertain{identifié} à » et « \uncertain{identifiés} » ; devant
$a \boxtimes b$, un terme raturé ; devant le $\underline{\mathrm{Hom}}$ de
la première ligne, un premier « Hom » biffé. Sous la seconde ligne, un signe $\simeq$
amorcé, sans suite}

\[
{}^{g}X(a) = p_{B*}\, \underline{\mathrm{Hom}}_{A \times B}({}^{g}a, X)
\qquad
{}^{d}X(b) = p_{B*}\, \underline{\mathrm{Hom}}_{A \times B}({}^{d}b, X)
\]
\note{l'indice du premier $p$ est surchargé ; on lit $B$ dans les deux
formules, alors que plus bas ${}^{g}X(L)$ s'écrit avec $p_{A*}$}

\struck{$\underline{\mathrm{Hom}}(L \boxtimes M, X) = X(L, M)$}
\[
\begin{aligned}
X(L, M) &= \underline{\mathrm{Hom}}_{A \times B}(L \boxtimes M, X)
 = \Gamma_{A \times B}\, \underline{\mathrm{Hom}}(L \boxtimes M, X)\\
{}^{g}X(L) &= p_{A*}\, \underline{\mathrm{Hom}}_{A \times B}({}^{g}L, X)
 = \underline{\mathrm{Hom}}_B({}^{g}L, X)\\
{}^{d}X(M) &= \underbrace{p_{B*}\, \underline{\mathrm{Hom}}_{A \times B}({}^{d}M, X)}_{X^{M}}
 = \underline{\mathrm{Hom}}_A({}^{d}M, X)
\end{aligned}
\]
\[
\underline{\mathrm{Hom}}_B(X, Y^{M}) \simeq \underline{\mathrm{Hom}}_B(X \otimes M, Y)
\]

\page{13}

\note{en haut, au milieu du bord perforé, la lettre « f »}

\[
\begin{aligned}
(A \times B)^{\wedge} &\overset{\text{déf}}{=}
\underline{\mathrm{Hom}}\bigl((A \times B)^{\circ}, (\mathrm{Ens})\bigr)
\simeq \underline{\mathrm{Hom}}\bigl(A^{\circ} \times B^{\circ}, (\mathrm{Ens})\bigr)\\
&\simeq \underline{\mathrm{Hom}}(A^{\circ}, B^{\wedge})
\simeq \underline{\mathrm{Hom}}_{!}(A^{\wedge\circ}, B^{\wedge})\\
&\simeq \underline{\mathrm{Hom}}(B^{\circ}, A^{\wedge})
\simeq \underline{\mathrm{Hom}}_{!}(B^{\wedge\circ}, A^{\wedge})\\
&\simeq \underline{\mathrm{Hom}}_{!!}\bigl((A^{\wedge})^{\circ} \times (B^{\wedge})^{\circ}, \mathrm{Ens}\bigr)
\simeq \underline{\mathrm{Hom}}_{!}\bigl((A^{\wedge} \times B^{\wedge})^{\circ}, \mathrm{Ens}\bigr)
\end{aligned}
\]
\note{dans la deuxième ligne, $B^{\wedge}$ est récrit sur un premier
$(\mathrm{Ens})$ et suivi d'un signe raturé}

\marginal{le signe ! signifie : foncteurs commutant aux $\varprojlim$
quelconques ! le signe !! : et le même, \uncertain{contigus} \ill{}
\struck{\ill{}} des \uncertain{variables} \uncertain{séparément}
\uncertain{comme} \ill{} \ill{} et un seul signe ! relatif :
$(A^{\wedge} \times B^{\wedge})^{\circ}$}

Si $X \in (A \times B)^{\wedge}$, le foncteur \struck{\ill{}}
correspondant
\[
(L, M) \longmapsto X(L, M) \quad :
\underset{(A^{\wedge} \times B^{\wedge})^{\circ}}{\underbrace{A^{\wedge\circ} \times B^{\wedge\circ}}}
\longrightarrow (\mathrm{Ens})
\]
est donné par
\[
X(L, M) = \mathrm{Hom}(L \boxtimes M, X) \quad \text{où} \quad
L \boxtimes M = p_B^{*}(L) \times p_A^{*}(M) .
\]
\note{après le premier membre, un « $= \mathrm{Hom}$ » biffé ; le signe
entre $p_B^{*}(L)$ et $p_A^{*}(M)$ est surchargé, un $\boxtimes$ récrit en
produit. Au début de la ligne qui précède, dans $(L, M)$, le $L$ est
récrit sur un $X$}

\textbf{\uncertain{Proposition}}\note{« Proposition », souligné, est écrit
au-dessus de \struck{Corollaire}, biffé} Le foncteur $(L, M) \longmapsto
L \boxtimes M$ de $(A^{\wedge} \times B^{\wedge})^{\circ}$ dans
$(A^{\wedge} \times B^{\wedge})^{\wedge}$ est pleinement fidèle, il
s'identifie au foncteur qui associe à tout objet $(L, M)$ de
$A^{\wedge} \times B^{\wedge}$ le foncteur
\struck{\uncertain{Can.} $(A^{\wedge} \times B^{\wedge})^{\circ}
\to \underline{\mathrm{Hom}}\bigl((A^{\wedge} \times B^{\wedge})^{\circ},$}
\struck{\uncertain{contravariant} qu'il représente} sur
$A^{\wedge} \times B^{\wedge}$ qu'il représente, i.e.
\note{à la fin de la première ligne, au-dessus de $(A \times B)^{\wedge}$
raturé, est récrit $(A^{\wedge} \times B^{\wedge})^{\wedge}$ ; au-dessus de
« qui associe », « tout » et « objet » sont ajoutés}
\[
(L \boxtimes M)(L', M') \simeq
\mathrm{Hom}_{A^{\wedge} \times B^{\wedge}}\bigl((L', M'), (L, M)\bigr)
= \mathrm{Hom}(L', L) \times \mathrm{Hom}(M', M)
\]
\[
\overset{\text{déf}}{\|} \qquad
\mathrm{Hom}(L' \boxtimes M', L \boxtimes M)
\]
\note{le dernier terme de la première ligne est serré contre le bord de
la feuille et surchargé ; « $(M', M)$ » est une lecture douteuse}

\struck{Bien sûr, par fonctorialité on a}
\struck{$\mathrm{Hom}(L', L) \times \mathrm{Hom}(M', M) \longrightarrow
\mathrm{Hom}(L' \boxtimes M', L \boxtimes M)$}

En effet, la formule est vraie par définition de $L \boxtimes M$ si
\struck{\ill{}} $L' = a \in \mathrm{Ob}\, A$, $M' = b \in \mathrm{Ob}\, B$,
\struck{et \ill{}} D'autre part, pour $L, M$ fixés, et $L', M'$ variables,
les deux membres, comme foncteurs en $L', M'$, transforment
$\varinjlim$ en $\varprojlim$ \struck{comme il \ill{} \ill{}} \struck{et
il \uncertain{prolonge}} dans les deux variables.\note{les signes
$\varinjlim$ et $\varprojlim$ sont écrits « lim » souligné d'une flèche ;
les deux lignes biffées qui précèdent la fin de la phrase sont de lecture
très douteuse}

\textbf{Corollaire}\quad
$\underline{\mathrm{Hom}}(L' \boxtimes M', L \boxtimes M) \simeq
\underline{\mathrm{Hom}}(L', L) \boxtimes \underline{\mathrm{Hom}}(M', M)$
\note{devant $L \boxtimes M$, une lettre raturée à l'encre épaisse}

\emph{Dém.} Il suffit de prendre les $\mathrm{Hom}(a \boxtimes b, -)$ pour
$(a, b) \in \mathrm{Ob}\, A \times B$
\[
\begin{array}{c}
\mathrm{Hom}_{A \times B}\bigl(\underbrace{(a \boxtimes b) \times (L' \boxtimes M')}_{(a \times L') \boxtimes (b \times M')}, L \boxtimes M\bigr)\\
\|\\
\mathrm{Hom}_A(a \times L', L) \times \mathrm{Hom}_B(b \times M', M)
\end{array}
\]
\[
\begin{array}{c}
\mathrm{Hom}_A\bigl(a, \underline{\mathrm{Hom}}(L', L)\bigr) \times \mathrm{Hom}_B\bigl(b, \underline{\mathrm{Hom}}(M', M)\bigr)\\
\|\\
\mathrm{Hom}_A(a \times L', L) \times \mathrm{Hom}_B(b \times M', M)
\end{array}
\]
\note{sur la page, les deux colonnes sont côte à côte ; on les donne l'une sous l'autre. Colonne de gauche : un premier argument $a \times L' \boxtimes M'$
est biffé, et récrit au-dessus $(a \boxtimes b) \times (L' \boxtimes M')$,
avec l'accolade en dessous. Un signe $\|$ en tête de colonne relie le
tout à la ligne « Corollaire » ; la colonne de droite descend de même de
$\underline{\mathrm{Hom}}(L', L) \boxtimes
\underline{\mathrm{Hom}}(M', M)$. Au pied de la colonne de droite,
« ok »}

\page{14}

\note{écrit sur la moitié libre d'une feuille de listing ; pas de lettre
de page. Formules isolées, sans texte suivi}

\[
\begin{aligned}
(X^{\Delta_n})_0 &= p_{B*}\bigl(\underline{\mathrm{Hom}}(p_A^{*}(\Delta_n), X)\bigr) = X_n\\
(X^{\dot{\Delta}_n})_0 &= p_{B*}\, \underline{\mathrm{Hom}}(\dot{\Delta}_n, X) = \varprojlim\,(X_{n-1}, X_{n-2}\ \ldots)
\end{aligned}
\]
\note{dans la première ligne, devant $p_{B*}$, un « $\underline{\mathrm{Hom}}_B$ » biffé ; le $X$ de $X_n$ est surchargé ; dans la
seconde, les arguments de la limite sont de lecture douteuse}

\[
X^{\Delta_n} \longrightarrow X^{\dot{\Delta}_n} \underset{Y^{\dot{\Delta}_n}}{\times} Y^{\Delta_n}
\qquad\qquad
X_n \longrightarrow (X^{\dot{\Delta}_n})_0 \underset{(Y^{\dot{\Delta}_n})_0}{\times} Y_n
\]

\note{à droite, $(\mathrm{Cosk}_{n-1} X)_n$, suivi d'un terme biffé ; un
petit triangle hachuré dont les trois côtés sont marqués $x_1$ ; plus bas,
« $A_n = A$ » au-dessus d'un « $A$ », et trois traits parallèles joignant
$X_{n-2}$ à $X_{n-1}$}

\[
\underline{\mathrm{Hom}}(\dot{\Delta}_n, a)
\qquad
\underline{\mathrm{Hom}}(M, a) \overset{??}{\Longrightarrow} a
\]
\[
\mathrm{Hom}(X \times M, a) \simeq \mathrm{Hom}(X, -)
\]
\[
\mathrm{Hom}\bigl(\alpha \boxtimes (\beta \times M), \alpha \boxtimes e_B\bigr)
\simeq \mathrm{Hom}(\alpha \boxtimes \beta, \alpha \boxtimes e_B)
\]
\note{la dernière formule touche le bord de la feuille ; devant le second $\alpha$ du membre de gauche, une lettre biffée ; le second
argument du membre de droite est de lecture douteuse. Dessous, un
« Hom » et un « $L$ » isolés, sans suite}

\page{15}

\note{en haut, au milieu du bord perforé, la lettre « g »}

\textbf{Corollaire}\quad $p_B^{*} : A^{\wedge} \to (A \times B)^{\wedge}$
commute aux $\underline{\mathrm{Hom}}$ (et \uncertain{de} symétrique
pour $p_A^{*}$).

Il suffit de prendre \struck{$M' \simeq e_{B^{\wedge}}$} $M =$
\uncertain{$M' = M = e_{B^{\wedge}}$}. On trouve \ill{}
\[
\underline{\mathrm{Hom}}(L' \boxtimes M', L \boxtimes e_{B^{\wedge}})
\simeq \underline{\mathrm{Hom}}(L', L) \boxtimes e_{B^{\wedge}}
\]
\note{devant $L \boxtimes e_{B^{\wedge}}$, une lettre surchargée}

En particulier, faisant $L' = e_{A^{\wedge}}$ et changeant $M'$ en $M$
\[
\underline{\mathrm{Hom}}(e_{A^{\wedge}} \boxtimes M, L \boxtimes e_{B^{\wedge}})
\simeq L \boxtimes e_{B^{\wedge}}
\]
\note{sous le $\simeq$, une petite flèche $\leftarrow$ ; dans « $L' =
e_{A^{\wedge}}$ », le $e$ est récrit sur une lettre raturée}

--- ce qui \uncertain{correspond} à la \struck{\ill{}}
\uncertain{flèche} \ill{} flèche canonique
\[
\underline{\mathrm{Hom}}(P, Q) \longleftarrow Q
\]
valable dans \uncertain{toute} catégorie, \struck{\ill{} produits}
\struck{\ill{}}
\note{deux mots au-dessus de la ligne sont de lecture douteuse ;
sous « catégorie », deux lignes biffées, la première lue
\uncertain{produits}}

Dans la notation $X^{M}$, si on identifie $M$ à $e_{A^{\wedge}} \boxtimes M
= p_A^{*}(M)$, et $L$ à $L \boxtimes e_{B^{\wedge}} \simeq p_B^{*}(L)$, on
trouve donc
\[
L^{M} = L \qquad \text{pour } L \in A^{\wedge},\ M \in \mathrm{Ob}\, B^{\wedge}
\]
\note{dans « $L \boxtimes e_{B^{\wedge}}$ », le $L$ est récrit sur une
autre lettre ; dans « $L \in A^{\wedge}$ », un « Ob » est raturé}

\page{16}

\note{écrit sur la moitié libre d'une feuille de listing ; pas de lettre
de page. Un long trait oblique traverse toute la page, du haut au milieu
jusqu'au listing : la page semble annulée ; elle est transcrite telle
qu'elle se lit. Ici $M$ désigne une catégorie, non un préfaisceau}

\[
\underline{\mathrm{Hom}}(A, M) \simeq \underline{\mathrm{Hom}}_{!}(A^{\wedge}, M)
\]
\note{sous $M$ : « \uncertain{stable} \uncertain{par} $\varinjlim$ »,
écrit sur deux lignes. À droite : « ! = \struck{\ill{}} commutant aux
$\varinjlim$ », « lim » souligné d'une flèche vers la droite}

\[
\begin{aligned}
\underline{\mathrm{Hom}}(A^{\circ}, M) &\simeq \underline{\mathrm{Hom}}_{!}(A^{\circ\wedge}, M)
\simeq \Bigl(\underline{\mathrm{Hom}}^{!}\bigl(\underbrace{A^{\circ\wedge\circ} = A^{\vee\circ}}, M^{\circ}\bigr)\Bigr)^{\circ}\\
\wr\quad &\\
\underline{\mathrm{Hom}}(A, M^{\circ})^{\circ} &\simeq \bigl(\underline{\mathrm{Hom}}_{!}(A^{\wedge}, M^{\circ})\bigr)^{\circ}
= \underline{\mathrm{Hom}}^{!}(A^{\wedge\circ}, M)
\end{aligned}
\]
\note{sous l'accolade : « \uncertain{catégorie} \uncertain{stable}
\uncertain{par} $\varprojlim$ » ; le $M^{\circ}$ de la première ligne est
récrit sur un $M$, et une parenthèse est raturée après lui}

\struck{$\underline{\mathrm{Hom}}^{!}(A^{\vee\circ}, M^{\circ}) \simeq
\underline{\mathrm{Hom}}^{!}_{!}(A^{\wedge}, M^{\circ})$}
\note{cette formule, encadrée, est barrée de plusieurs traits obliques ;
l'exposant du second $\underline{\mathrm{Hom}}$ est surchargé}

\[
\underline{\mathrm{Hom}}_{!}(A^{\vee}, M)^{\circ} \simeq \underline{\mathrm{Hom}}_{!}(A^{1}, M^{\circ})
\]
\[
\underline{\mathrm{Hom}}(A^{\circ}, M) \simeq \underline{\mathrm{Hom}}_{!}(A^{\circ\wedge}, M)
= \underline{\mathrm{Hom}}_{!}(A^{\vee}, M)
\]
\uncertain{donc}
\[
\underline{\mathrm{Hom}}(A, M^{\circ}) \simeq \underline{\mathrm{Hom}}_{!}(A^{\wedge}, M^{\circ})
\]
\note{dans la quatrième formule, l'exposant de $A$ est lu, avec doute,
« 1 »}

\section{Question des $\varinjlim$ filtrantes dans Hot}

\note{titre de sa main, seul en haut d'une feuille vierge (page 17) :
« Question des lim filtrantes dans Hot », « lim » souligné. Il annonce les
pages en anglais qui suivent, et celles du lot 2}

\page{18}

\note{en haut, cerclé, « a » : sa lettre de page. Cette suite en anglais
(a, b ici ; puis, au lot 2, e à la page 22 et d à la page 24) est une autre
suite que celle, en français, des pages 3 à 15}

Filtering \uncertain{limits} exist in $(\mathrm{Hot})$?

We would like moreover that \struck{the} the canonical functors
\[
(\mathrm{Cat}) \longrightarrow \mathrm{Hot}(\underline{W}), \qquad
A^{\wedge} \longrightarrow \mathrm{Hot}_A \simeq \mathrm{Hot}(\underline{W})
\quad (A \text{ a test category})
\]
should commute with such limits. The first case is reduced to the second
with $A = \Delta$, using the standard test-functor $\Delta
\xrightarrow{\;i\;} (\mathrm{Cat})$, giving rise to the Nerve-functor
$i^{*} : (\mathrm{Cat}) \to \Delta^{\wedge}$, as the latter commutes with
filtering direct limits (due to the fact that $i(a)$ is ``of f.p.'' for
any $a \in \Delta$.\note{l'indice de la flèche $i$ est surchargé ; la
parenthèse n'est pas refermée}

Take first the case of a countable direct system, we may assume it is
indexed by $\mathbb{N}$
\[
X_0 \xrightarrow{\;f_0\;} X_1 \xrightarrow{\;f_1\;} X_2 \cdots
\]
Using factorisation of any morphism as $pi$, with $i$ mono and $p$ a
homotopism, we get a diagram

\begin{tikzcd}[column sep=large]
X'_0 \arrow[r, hook, "f'_0"] \arrow[d, "{u_0 \,=}"'] & X'_1 \arrow[r, hook, "f'_1"] \arrow[d, "u_1"] & X'_2 \arrow[r, hook, "f'_2"] \arrow[d, "u_2"] & \cdots & X'_\infty \arrow[d, "u_\infty"] \\
X_0 \arrow[r, "f_0"'] & X_1 \arrow[r, "f_1"'] & X_2 \arrow[r, "f_2"'] & \cdots & X_\infty
\end{tikzcd}

where the $f'_i$ are mono, and the vertical arrows are homotopy
equivalences. Let $X'_\infty$, $X_\infty$ be the limits, hence
\[
u_\infty : X'_\infty \longrightarrow X_\infty .
\]
We've to check

(i) $X'_\infty$ \struck{\ill{}} a direct limit of the $X'_i$'s in
$(\mathrm{Hot}_A)$

(ii) $X'_\infty \to X_\infty$ is an iso in $(\mathrm{Hot})_A$, i.e.\ is
\struck{\ill{}} $\in W$.

\struck{(and equivalent, in fact)}\note{cette parenthèse, sous (ii), est
biffée ; elle est rattachée à la ligne suivante}

Point (i) is enough \struck{\ill{}} \uncertain{assuming} that
$(\mathrm{Hot}_A)$ is stable under countable filtering direct limits, and
that $A^{\wedge} \to (\mathrm{Hot}_A)$ commutes \struck{\ill{}} under
\uncertain{to} such limits, when the transition morphisms are mono.
\struck{\ill{}} As for (ii), granting (i), it \struck{is equivalent to}
\uncertain{is} with the desired result of commutation of $A^{\wedge} \to
(\mathrm{Hot}_A)$ to countable filtering direct limit.\note{« filtering »
est ajouté au-dessus de « direct » ; la fin de la phrase « it … with » est
de lecture douteuse, un « is equivalent to » biffé étant prolongé d'une
flèche vers la ligne suivante}

To check the limiting property
\[
(*) \qquad \mathrm{Hom}_{\mathrm{Hot}}(X_\infty, Y) \xrightarrow{\;\sim\;}
\varprojlim_i \mathrm{Hom}_{\mathrm{Hot}}(X_i, Y) ,
\]
\note{le $Y$ du premier membre est récrit sur une autre lettre}

\page{20}

\note{en haut, cerclé, « b » : sa lettre de page ; la page fait suite
à la page 18}

we may assume $Y$ is a $(\underline{W})$-Kan object of $A^{\wedge}$, in
which case
\[
\mathrm{Hom}_{\mathrm{Hot}}(X, Y) \simeq \overline{\mathrm{Hom}}(X, Y)
\]
(homotopy classes of maps). Moreover, because of the assumption on $Y$,
the homotopy relation \struck{of} \uncertain{among} maps $X \to Y$ is
\struck{equivalent to} the same as elementary homotopy.\note{« the »,
au-dessus de « classes », et « among », « the same », au-dessus des mots
biffés, sont ajoutés entre les lignes}

Injectivity of $(*)$? We get maps $g_0, g_1, \ldots$

\begin{tikzcd}
X_0 \arrow[r, hook] \arrow[dr, "g_0"'] & X_1 \arrow[r, hook] \arrow[d, "g_1"] & X_2 \arrow[r, hook] \arrow[dl, "g_2"] & \cdots & X_n \arrow[r, "f_n"] & X_{n+1} \\
 & Y & & & &
\end{tikzcd}

\note{sous $X_n$ et $X_{n+1}$, deux flèches marquées $g_n$ et $g_{n+1}$,
tracées vers la gauche sans atteindre $Y$}

such that $g_{n+1} f_n \sim g_n$ for any $n$, the question is to find
\struck{a representative} for each $n$ a $g'_n \sim g_n$ such that
\struck{equality} $g'_{n+1} f_n = g'_n$ holds. This brings us to the
slightly more general situation (where $i : X \to X'$ is
$f_n : X_n \to X_{n+1}$, where $p : Y \to Y'$ is the map $Y \to e$, and
$g = g_n$, $\varphi = g_{n+1}$)\note{dans la parenthèse, un mot biffé
après $Y \to$ ; « $f_n : X_n \to X_{n+1}$ » est une restitution douteuse
de ce que la ligne abrège}

\begin{tikzcd}[column sep=large, row sep=large]
X \arrow[r, "g"] \arrow[d, "i"'] & Y \arrow[d, "p"] \\
X' \arrow[r, "g'"'] \arrow[ur, dashed, "\varphi"] & Y'
\end{tikzcd}

\note{la diagonale $\varphi$ est tracée en tirets, sa lettre surchargée}

when $i$ is a cofibration, $p$ is a fibration, and the two triangles
commutative up to homotopy
\[
\varphi i \overset{k}{\sim} g \qquad p\varphi \overset{h}{\sim} g'
\]
(we assume $X, X'$ cofibrant, $Y, Y'$ fibrant, so that
$\overset{\ell}{\sim}$ and $\overset{r}{\sim}$ are the same, --- we
look at the situation in any model category)\note{les deux signes sont
$\sim$ surmontés d'une petite lettre, lus $\ell$ et $r$ (homotopies à
gauche et à droite) ; lecture douteuse}.

May we find a $\overline{\varphi} : X' \to Y$ (homotopic to $\varphi$,)
s.th.\ the two triangles commute:
\[
\overline{\varphi} \sim \varphi , \qquad \varphi i = g, \quad p\varphi = g' \ ?
\]
\note{« homotopic to $\varphi$, » est ajouté au-dessus de la ligne ; dans
la dernière ligne il écrit $\varphi$, non $\overline{\varphi}$}

\struck{Let $\overline{X}'$ be a cylinder object over $X'$}

Now this is about exactly the content of Quillen's version of ``covering
homotopy extension theorems'' for simplicial closed model categories
(prop.~4 of Chap.~II, $\mathrm{n}^{\circ}$~2.4) --- although it was a
little hard to make a sense of it, the way he stated it

\marginal{There is such a possibility, iff we can find $k$ and $h$
s.th.\ the two \uncertain{resulting} homotopies
$h \circ (i \otimes J)$, \ill{} : $p\varphi i \longrightarrow
p\overline{\varphi}\,i = g' i$ \uncertain{are} homotopic.
\uncertain{This} is presumably not always \uncertain{satisfied}!}
\note{note de sa main, écrite en travers dans la marge gauche ; devant,
un mot biffé, \struck{\ill{}}. Les lectures en sont très douteuses}


\end{document}
